\section{关系}
\subsection{无序对，有序对}
\begin{definition}
    由两个元素\(x_1\)和\(x_2\)组成的集合
    \begin{equation*}
        \Set{x_1, x_2}
    \end{equation*}
    称作\DefineConcept{无序对}（unordered pair）.
\end{definition}

\begin{example}
    %@see: 《Elements of Set Theory》 P38 Exercise 1.
    定义：
    \begin{equation*}
        \opair{x,y,z}^* = \{\{x\},\{x,y\},\{x,y,z\}\}.
    \end{equation*}
    找出元素\(u,v,w,x,y,z\)使
    \begin{equation*}
        \opair{x,y,z}^* = \opair{u,v,w}^*
    \end{equation*}
    成立，
    但\(y \neq v\)或\(z \neq w\).
    \begin{solution}
        取\(x=u=1\)、\(y=w=1\)、\(z=v=2\).
        集合中的重复元素不影响集合本身，
        因此
        \begin{align*}
            \opair{1,1,2}^*
            &= \Set{\Set{1},\Set{1,1},\Set{1,1,2}}
            = \Set{\Set{1},\Set{1,2}}, \\
            \opair{1,2,1}^*
            &= \Set{\Set{1},\Set{1,2},\Set{1,2,1}}
            = \Set{\Set{1},\Set{1,2}}.
        \end{align*}
        两个集合相等，
        但\(y=1\neq2=v\)，
        且\(z=2\neq1=w\)，
        故满足题设条件.
    \end{solution}
\end{example}

\begin{definition}
    两个元素\(x_1\)和\(x_2\)按一定顺序形成的排列称为\DefineConcept{有序对}，
    记作
    \begin{equation*}
        \opair{x_1,x_2}.
    \end{equation*}

    我们可以采用多种方式构造集合用于表记有序对\(\opair{x_1,x_2}\)，
    例如（由Norbert Wiener于1914年构造的）
    \begin{equation*}
        \Set{ \Set{ \Set{x_1}, \emptyset }, \Set{ \Set{ x_2 } } }.
    \end{equation*}
    但我们更常用（由Kazimierz Kuratowski于1921年构造的）以下形式
    \begin{equation*}
        \Set{ \Set{x_1}, \Set{x_1, x_2} }.
    \end{equation*}
    以后我们都采用第二种形式表记有序对，因此
    \begin{equation*}
        \opair{x_1,x_2}
        \defeq
        \Set{ \Set{x_1}, \Set{x_1, x_2} }.
    \end{equation*}
\end{definition}

有序对具有以下性质.
\begin{property}\label{theorem:集合论.有序对的性质1}
    %@see: 《Elements of Set Theory》 P36 Theorem 3A
    \(\opair{u,v} = \opair{x,y}\)的充分必要条件是：\(u=x\)且\(v=y\).
    \begin{proof}
        充分性.
        若\(u=x\)且\(v=y\)，
        代入有序对的定义便得\(\opair{u,v}=\opair{x,y}\).

        必要性.
        设\(\opair{u,v}=\opair{x,y}\).
        有序对作为集合非空，
        对等式两端分别取交集、并集可得
        \begin{align*}
            \Set{u}
            &= \bigcap\opair{u,v}
            = \bigcap\opair{x,y}
            = \Set{x}, \\
            \Set{u,v}
            &= \bigcup\opair{u,v}
            = \bigcup\opair{x,y}
            = \Set{x,y}.
        \end{align*}
        第一个等式给出\(u=x\).
        若\(v=u\)，
        则第二个等式成为\(\Set{x}=\Set{x,y}\)，
        从而\(y=x=v\).
        若\(v\neq u\)，
        则\(v\in\Set{x,y}\)且\(v\neq x\)，
        从而\(v=y\).
        因此总有\(u=x\)且\(v=y\).
    \end{proof}
\end{property}
\cref{theorem:集合论.有序对的性质1} 让我们可以不含糊地将有序对\(\opair{x,y}\)
中的\(x\)和\(y\)分别称为%
有序对的\DefineConcept{第一坐标}（first coordinate）%
和\DefineConcept{第二坐标}（second coordinate）.

\begin{example}
    设\(w\)是集合.
    为以下两个类选取记号，
    说明它们的含义以及它们为何是集合：
    \begin{align*}
        &\Set{x\given(\exists u)(\exists v)[w=\opair{u,v}\land x\in u]}, \\
        &\Set{y\given(\exists u)(\exists v)[w=\opair{u,v}\land y\in v]}.
    \end{align*}
    \begin{solution}
        在本例中分别将它们记作\(C_1(w)\)、\(C_2(w)\).
        若\(w=\opair{u,v}\)，
        则由有序对坐标的唯一性，
        对任意\(t\)有
        \begin{align*}
            t\in C_1(w)&\iff t\in u, \\
            t\in C_2(w)&\iff t\in v.
        \end{align*}
        因此\(C_1(w)=u\)、\(C_2(w)=v\)，
        即它们分别给出\(w\)的第一、第二坐标本身，
        而不是\(\Set{u}\)、\(\Set{v}\).
        若\(w\)不是有序对，
        则两个定义中的存在条件均不成立，
        所以\(C_1(w)=C_2(w)=\emptyset\).

        还可直接给出统一的集合上界.
        若\(w=\opair{u,v}\)，
        则\(\bigcup w=\Set{u,v}\)，
        从而\(\bigcup\bigcup w=u\cup v\).
        因此原来两个类分别等于
        \begin{align*}
            C_1(w)
            &=\Set*{t\in\bigcup\bigcup w
            \given(\exists u)(\exists v)[w=\opair{u,v}\land t\in u]}, \\
            C_2(w)
            &=\Set*{t\in\bigcup\bigcup w
            \given(\exists u)(\exists v)[w=\opair{u,v}\land t\in v]}.
        \end{align*}
        并集公理和子集公理保证它们都是集合.
        这里的记号只是在有序对之外延拓为空集的坐标提取记号，
        不表示非有序对也具有第一、第二坐标.
    \end{solution}
\end{example}

\subsection{直积}
\begin{definition}\label{definition:集合论.直积}
    设\(A,B\)都是集合.
    在\(A\)中任取一个元素\(x\)，
    在\(B\)中任取一个元素\(y\)，
    组成一个有序对\(\opair{x,y}\)，
    把这样的有序对作为新的元素，
    它们全体组成的类称为
    “\(A\)与\(B\)的\DefineConcept{直积}（direct product）”，
    或“\(A\)与\(B\)的\DefineConcept{笛卡尔乘积}（Cartesian product）”，
    记为\(A \times B\)，
    即
    \begin{equation}
        A \times B
        \defeq
        \Set*{ \opair{x,y} \given x \in A \land y \in B }.
    \end{equation}
    %@see: https://mathworld.wolfram.com/DirectProduct.html
\end{definition}

在宣称\cref{definition:集合论.直积} 合法以前，
我们必须确认它给出的类\(A \times B\)是不是真的是一个集合.
而要证明\(A \times B\)是一个集合，
我们就必须找到一个足够大的集合，
它包括了我们需要的全部有序对，
然后利用子集公理证明\(A \times B\)是这个集合的子集.
为此我们给出下述引理.
\begin{lemma}\label{theorem:集合论.有序对是其坐标元素所在集合的二重幂集的元素}
    %@see: 《Elements of Set Theory》 P37 Lemma 3B
    如果\(x,y \in A\)，那么\(\opair{x,y} \in \Powerset\Powerset A\).
    \begin{proof}
        由\(x,y\in A\)可知\(\Set{x},\Set{x,y}\subseteq A\)，
        即\(\Set{x},\Set{x,y}\in\Powerset A\).
        因此
        \begin{align*}
            \opair{x,y}
            &= \Set{\Set{x},\Set{x,y}}
            \subseteq\Powerset A.
        \end{align*}
        根据幂集的定义，
        \(\opair{x,y}\in\Powerset\Powerset A\).
    \end{proof}
\end{lemma}

\begin{theorem}\label{theorem:集合论.直积存在定理}
    %@see: 《Elements of Set Theory》 P38 Corollary 3C
    设\(A,B\)都是集合，那么存在这样一个集合\(C\)，使得\(C = A \times B\).
    \begin{proof}
        由配对公理、并集公理和幂集公理，
        \(\Powerset\Powerset(A\cup B)\)是集合.
        利用子集公理构造
        \begin{align*}
            C
            &= \Set*{w\in\Powerset\Powerset(A\cup B)
            \given (\exists x\in A)(\exists y\in B)[w=\opair{x,y}]}.
        \end{align*}
        若\(w\in C\)，
        则存在\(x\in A\)、\(y\in B\)使\(w=\opair{x,y}\)，
        因而\(w\in A\times B\).
        反之，
        若\(w\in A\times B\)，
        则存在这样的\(x,y\).
        由于\(x,y\in A\cup B\)，
        由前一引理可知\(w\in\Powerset\Powerset(A\cup B)\)，
        从而\(w\in C\).
        由外延公理得到\(C=A\times B\).
        此构造也适用于\(A\)或\(B\)为空集的情形.
    \end{proof}
\end{theorem}
\cref{theorem:集合论.直积存在定理} 告诉我们，
任意两个集合\(A\)和\(B\)的直积\(A \times B\)也是集合.

容易发现，对于任意集合\(A\)，总有
\begin{equation}
    A \times \emptyset
    = \emptyset \times A
    = \emptyset \times \emptyset
    = \emptyset.
\end{equation}
这是因为\((\forall x)[x \notin \emptyset]\).

\begin{example}
    设\(A,B,C\)是集合.
    证明：
    \begin{equation}
        A \subseteq B
        \implies
        A \times C \subseteq B \times C.
    \end{equation}
    \begin{proof}
        设\(A\subseteq B\)，
        任取\(w\in A\times C\).
        由直积的定义，
        存在\(x\in A\)、\(y\in C\)使\(w=\opair{x,y}\).
        由\(A\subseteq B\)可知\(x\in B\)，
        因而\(w\in B\times C\).
        故\(A\times C\subseteq B\times C\).
        若\(A\times C=\emptyset\)，
        则上述包含关系同样成立.
    \end{proof}
\end{example}

\begin{example}
    %@see: 《Elements of Set Theory》 P38 Exercise 2.(a)
    %@see: 《Elements of Set Theory》 P65 Exercise 54.
    设\(A,B,C\)是集合.
    证明：
    \begin{gather}
        A \times (B \cup C) = (A \times B) \cup (A \times C),
        \label{equation:集合论.直积分配律1} \\
        A \times (B \cap C) = (A \times B) \cap (A \times C),
        \label{equation:集合论.直积分配律2} \\
        A \times (B - C) = (A \times B) - (A \times C).
        \label{equation:集合论.直积分配律3}
    \end{gather}
    \begin{proof}
        三个等式两端的元素都是有序对，
        因此只需对任意\(x,y\)检验\(\opair{x,y}\)的归属.
        对于并集，
        有
        \begin{align*}
            \opair{x,y}\in A\times(B\cup C)
            &\iff x\in A\land(y\in B\lor y\in C) \\
            &\iff (x\in A\land y\in B)\lor(x\in A\land y\in C) \\
            &\iff\opair{x,y}\in(A\times B)\cup(A\times C).
        \end{align*}
        对于交集，
        有
        \begin{align*}
            \opair{x,y}\in A\times(B\cap C)
            &\iff x\in A\land y\in B\land y\in C \\
            &\iff (x\in A\land y\in B)\land(x\in A\land y\in C) \\
            &\iff\opair{x,y}\in(A\times B)\cap(A\times C).
        \end{align*}
        对于差集，
        有
        \begin{align*}
            \opair{x,y}\in A\times(B-C)
            &\iff x\in A\land y\in B\land y\notin C \\
            &\iff (x\in A\land y\in B)\land\neg(x\in A\land y\in C) \\
            &\iff\opair{x,y}\in(A\times B)-(A\times C).
        \end{align*}
        最后一组推导的第二步利用了\(x\in A\).
        由外延公理，
        三个等式均成立.
    \end{proof}
\end{example}

\begin{example}
    %@see: 《Elements of Set Theory》 P38 Exercise 2.(b)
    证明：如果\(A \times B = A \times C\)且\(A \neq \emptyset\)，则\(B = C\).
    \begin{proof}
        由于\(A\neq\emptyset\)，
        可以固定一个\(a\in A\).
        对任意集合\(y\)，
        由有序对坐标的唯一性和题设等式可得
        \begin{align*}
            y\in B
            &\iff\opair{a,y}\in A\times B \\
            &\iff\opair{a,y}\in A\times C \\
            &\iff y\in C.
        \end{align*}
        由外延公理，
        \(B=C\).
        非空条件使得上述\(a\)可以选取；
        当\(A=\emptyset\)时，
        两个直积无论\(B,C\)为何均相等，
        因而不能由此推出\(B=C\).
    \end{proof}
\end{example}

\begin{example}
    %@see: 《Elements of Set Theory》 P38 Exercise 3.
    证明：\(A \times \bigcup B = \bigcup\Set{ A \times X \given X \in B }\).
    \begin{proof}
        先说明右端的集族是集合.
        对每个\(X\in B\)，
        有\(X\subseteq\bigcup B\)，
        因而\(A\times X\subseteq A\times\bigcup B\).
        利用幂集公理和子集公理构造
        \begin{align*}
            \mathcal{F}
            &=\Set*{T\in\Powerset(A\times\bigcup B)
            \given(\exists X\in B)[T=A\times X]}.
        \end{align*}
        它恰好是\(\Set{A\times X\given X\in B}\)，
        故并集公理保证\(\bigcup\mathcal{F}\)也是集合.

        对任意\(x,y\)，
        有
        \begin{align*}
            \opair{x,y}\in A\times\bigcup B
            &\iff x\in A\land(\exists X\in B)[y\in X] \\
            &\iff(\exists X\in B)[x\in A\land y\in X] \\
            &\iff(\exists X\in B)[\opair{x,y}\in A\times X] \\
            &\iff\opair{x,y}\in\bigcup\mathcal{F}.
        \end{align*}
        两端都只含有序对，
        由外延公理即得所证等式.
        当\(A=\emptyset\)或\(B=\emptyset\)时，
        两端都为空集，
        结论仍成立.
    \end{proof}
\end{example}

\begin{example}
    %@see: 《Elements of Set Theory》 P38 Exercise 5.(a)
    设\(A,B\)都是集合，试证：\(\Set{ \{x\} \times B \given x \in A }\)也是集合.
    \begin{proof}
        若\(x\in A\)，
        则\(\Set{x}\times B\subseteq A\times B\)，
        即\(\Set{x}\times B\in\Powerset(A\times B)\).
        因此利用子集公理构造
        \begin{align*}
            \mathcal{F}
            &= \Set*{T\in\Powerset(A\times B)
            \given(\exists x\in A)[T=\Set{x}\times B]}.
        \end{align*}
        由定义，
        \(\mathcal{F}=\Set{\Set{x}\times B\given x\in A}\)，
        因而题设集族是集合.
        这里构造的是由各个\(\Set{x}\times B\)组成的集族，
        而不是其中的某一个集合.
        若\(A=\emptyset\)，
        则\(\mathcal{F}=\emptyset\)；
        若\(A\neq\emptyset\)且\(B=\emptyset\)，
        则\(\mathcal{F}=\Set{\emptyset}\).
    \end{proof}
\end{example}

\subsection{关系}
\begin{definition}
    \DefineConcept{关系}（relation）是有序对的集合.
\end{definition}

如果有序对\(\opair{x,y}\)是关系\(\rel{R}\)的元素，即
\begin{equation*}
    \opair{x,y} \in \rel{R},
\end{equation*}
那么称“\(x\)与\(y\)有\(\rel{R}\)关系”，
记作\(x\rel{R}y\)；
反之，
\begin{equation*}
    \opair{x,y} \notin \rel{R},
\end{equation*}
那么称“\(x\)与\(y\)没有\(\rel{R}\)关系”.

给定集合\(X\)和集合\(Y\)，
如果关系\(\rel{R}\)满足
\begin{equation*}
    \rel{R} \subseteq X \times Y,
\end{equation*}
我们就称“\(\rel{R}\)是\(X\)与\(Y\)之间的\DefineConcept{二元关系}（binary relation）”.
特别地，当\(X = Y\)时，就称“\(\rel{R}\)是\(X\)上的\DefineConcept{二元关系}”.

从关系的定义可以看出，关系是有序对的集合.
即便是看起来毫无意义的，只由一个有序对组成的集合，也可以看成是一个关系.
我们自然会认为有的关系比别的关系更有意义，
例如我们马上会提到的“映射”“等价关系”和“排序关系”.

\begin{definition}\label{definition:集合论.定义域与值域的定义}
    设\(R\)是集合.

    如果集合\(A\)满足
    \begin{equation*}
        x \in A \iff (\exists y)[\opair{x,y} \in R],
    \end{equation*}
    那么称\(A\)为“\(R\)的\DefineConcept{定义域}（domain）”，
    记作\(\dom R\)，
    即
    \begin{equation*}
        \dom R
        \defeq
        \Set{
        x
        \given
        (\exists y)[\opair{x,y} \in R]
        }.
    \end{equation*}

    如果集合\(B\)满足
    \begin{equation*}
        x \in B \iff (\exists t)[\opair{t,x} \in R],
    \end{equation*}
    那么称\(B\)为“\(R\)的\DefineConcept{值域}（range）”，
    记作\(\ran R\)，
    即
    \begin{equation*}
        \ran R
        \defeq
        \Set{
        x
        \given
        (\exists t)[\opair{t,x} \in R]
        }.
    \end{equation*}

    最后，我们把\(R\)的定义域与它的值域的并集称为
    “\(R\)的\DefineConcept{域}（field）”，记作\(\fld R\)，即
    \begin{equation*}
        \fld R \defeq \dom R \cup \ran R.
    \end{equation*}
\end{definition}

为了正当化\cref{definition:集合论.定义域与值域的定义}，
我们必须明确：对于任意集合\(R\)，存在这样一个集合，
\(R\)中的有序对的第一坐标和第二坐标都是这个集合的元素.
这个问题类似于我们之前对于直积\(A \times B\)的定义的正当化，
当时我们证明了\cref{theorem:集合论.直积存在定理}.
为此，我们给出下述引理，
可以注意到它与\cref{theorem:集合论.有序对是其坐标元素所在集合的二重幂集的元素} 的关联.

\begin{lemma}
    %@see: 《Elements of Set Theory》 P41 Lemma 3D
    如果\(\opair{x,y} \in A\)，那么\(x,y \in \bigcup\bigcup A\).
    \begin{proof}
        根据有序对的定义，
        \(\Set{x,y}\in\opair{x,y}\).
        又由\(\opair{x,y}\in A\)，
        可得\(\Set{x,y}\in\bigcup A\).
        由于\(x,y\in\Set{x,y}\)，
        再次应用并集的定义便得\(x,y\in\bigcup\bigcup A\).
    \end{proof}
\end{lemma}
这条引理已经在\cref{example:集合论.有序对各坐标的取值范围} 得到证明.
利用这条引理，再加上子集公理，我们可以构造集合\(R\)的定义域和值域：
\begin{gather}
    \dom R
    \defeq
    \Set*{ x \in \bigcup\bigcup R \given (\exists y)[\opair{x,y} \in R] }, \\
    \ran R
    \defeq
    \Set*{ x \in \bigcup\bigcup R \given (\exists t)[\opair{t,x} \in R] }.
\end{gather}

\begin{example}
    %@see: 《Elements of Set Theory》 P41 Exercise 6.
    试证：集合\(A\)是一个关系的充分必要条件是
    \begin{equation*}
        A \subseteq \dom A \times \ran A.
    \end{equation*}
    \begin{proof}
        必要性.
        设\(A\)是关系，
        任取\(w\in A\).
        由关系的定义，
        存在\(x,y\)使\(w=\opair{x,y}\).
        由于\(\opair{x,y}\in A\)，
        根据定义域和值域的定义，
        \(x\in\dom A\)、\(y\in\ran A\).
        因而\(w\in\dom A\times\ran A\)，
        所以\(A\subseteq\dom A\times\ran A\).

        充分性.
        设\(A\subseteq\dom A\times\ran A\).
        直积的每个元素都是有序对，
        因而\(A\)的每个元素也是有序对.
        又因为\(A\)是集合，
        所以\(A\)是关系.
        这也包括\(A=\emptyset\)的情形.
    \end{proof}
\end{example}

\begin{example}
    %@see: 《Elements of Set Theory》 P41 Exercise 7.
    试证：给定关系\(\rel{R}\)，总有
    \begin{equation*}
        \fld \rel{R} = \bigcup\bigcup \rel{R}.
    \end{equation*}
    \begin{proof}
        先证\(\fld\rel{R}\subseteq\bigcup\bigcup\rel{R}\).
        若\(t\in\fld\rel{R}\)，
        则\(t\in\dom\rel{R}\)或\(t\in\ran\rel{R}\).
        因而存在\(\rel{R}\)中的有序对以\(t\)为某一坐标.
        由前面的坐标取值范围引理，
        \(t\in\bigcup\bigcup\rel{R}\).

        再证反向包含关系.
        若\(t\in\bigcup\bigcup\rel{R}\)，
        则存在\(s\in\bigcup\rel{R}\)使\(t\in s\)，
        并且存在\(p\in\rel{R}\)使\(s\in p\).
        由于\(\rel{R}\)是关系，
        可写成\(p=\opair{x,y}=\Set{\Set{x},\Set{x,y}}\).
        从而\(s=\Set{x}\)或\(s=\Set{x,y}\)，
        所以\(t=x\)或\(t=y\).
        又由\(\opair{x,y}\in\rel{R}\)得
        \(x\in\dom\rel{R}\)、\(y\in\ran\rel{R}\)，
        因而\(t\in\fld\rel{R}\).
        由外延公理，
        两个集合相等.
    \end{proof}
\end{example}

\begin{example}
    %@see: 《Elements of Set Theory》 P41 Exercise 8.
    试证：对于任意集合\(A\)，总有
    \begin{gather}
        \dom\bigcup A = \bigcup\Set{ \dom\rel{R} \given \rel{R} \in A }, \\
        \ran\bigcup A = \bigcup\Set{ \ran\rel{R} \given \rel{R} \in A }.
    \end{gather}
    \begin{proof}
        先说明右端的两个集族存在.
        令\(U=\bigcup\bigcup\bigcup A\).
        对任意\(\rel{R}\in A\)，
        有\(\rel{R}\subseteq\bigcup A\).
        因而\(\rel{R}\)中每个有序对的两个坐标都属于\(U\)，
        即\(\dom\rel{R},\ran\rel{R}\subseteq U\).
        利用幂集公理和子集公理构造
        \begin{align*}
            \mathcal{D}
            &=\Set*{T\in\Powerset U
            \given(\exists\rel{R}\in A)[T=\dom\rel{R}]}, \\
            \mathcal{E}
            &=\Set*{T\in\Powerset U
            \given(\exists\rel{R}\in A)[T=\ran\rel{R}]}.
        \end{align*}
        它们分别就是题设的定义域集族和值域集族.

        对任意\(z\)，
        按定义展开可得
        \begin{align*}
            z\in\dom\bigcup A
            &\iff(\exists y)[\opair{z,y}\in\bigcup A] \\
            &\iff(\exists\rel{R}\in A)(\exists y)[\opair{z,y}\in\rel{R}] \\
            &\iff(\exists\rel{R}\in A)[z\in\dom\rel{R}] \\
            &\iff z\in\bigcup\mathcal{D}.
        \end{align*}
        同样有
        \begin{align*}
            z\in\ran\bigcup A
            &\iff(\exists x)[\opair{x,z}\in\bigcup A] \\
            &\iff(\exists\rel{R}\in A)(\exists x)[\opair{x,z}\in\rel{R}] \\
            &\iff(\exists\rel{R}\in A)[z\in\ran\rel{R}] \\
            &\iff z\in\bigcup\mathcal{E}.
        \end{align*}
        分别应用外延公理即可得到两个等式.
        这里不要求\(A\)的元素都是关系，
        因为定义域和值域只取决于各集合中的有序对.
        若\(A=\emptyset\)，
        则\(\mathcal{D}=\mathcal{E}=\emptyset\)，
        两个等式仍成立.
    \end{proof}
\end{example}

\subsection{多元关系}
我们可以将“有序对”和“二元关系”的概念分别扩展为“元组”和“多元关系”.

例如定义
\begin{equation*}
    \opair{x_1,x_2,x_3}
    \defeq
    \opair{\opair{x_1,x_2},x_3},
\end{equation*}
称之为\DefineConcept{三元组}（triple）；
定义
\begin{equation*}
    \opair{x_1,x_2,x_3,x_4}
    \defeq
    \opair{\opair{x_1,x_2,x_3},x_4},
\end{equation*}
称之为\DefineConcept{四元组}（quadruple）；
定义
\begin{equation*}
    \opair{x_1,x_2,x_3,x_4,x_5}
    \defeq
    \opair{\opair{x_1,x_2,x_3,x_4},x_5},
\end{equation*}
称之为\DefineConcept{五元组}（quintuple）；
以此类推，对于任意给定\(n\)，可以定义\(n\)元组（\(n\)-tuple）：
\begin{equation*}
    \opair{x_1,x_2,\dotsc,x_n}
    \defeq
    \opair{\opair{x_1,x_2,\dotsc,x_{n-1}},x_n}.
\end{equation*}
为了让我们的定义看起来整齐划一，我们还可以补充定义
\begin{equation*}
    \opair{x} \defeq x,
\end{equation*}
称之为\DefineConcept{一元组}（1-tuple）.

我们把“在\(A\)上的\(n\)元\DefineConcept{关系}（n-ary \emph{relation} on \(A\)）”
定义为由\(n\)元组构成的集合.
由于\(A\)上的二元关系是\(A \times A\)的子集，
以及\(A\)上的三元关系（ternary relation, 3-ary relation）是\((A \times A) \times A\)的子集，
所以三元关系也可归结为一种二元关系；
同理，其他\(n\)元关系，只要\(n>1\)，也都可以归结为二元关系.
虽然在上面对\(n\)元关系的定义中，
我们也把由\(A\)中包括的一元组构成的集合称为
“\(A\)上的一元关系（unary relation, 1-ary relation）”，
但它只是\(A\)的一个子集，根本不满足关系的定义.

\subsection{投射}
\begin{definition}
    %@see: 《点集拓扑讲义（第四版）》（熊金城） P22 定义1.5.5
    设\(\AutoTuple{X}{n}\ (n\geq1)\)都是集合，
    \(X = X_1 \times \dotsb \times X_n\)是笛卡尔积.
    把
    \begin{equation*}
        \Set{
        \opair{
        \opair{\AutoTuple{x}{n}},
        x_i
        }
        \given
        (\AutoTuple{x}{n}) \in X
        }
        \quad(i=1,2,\dotsc,n)
    \end{equation*}
    称为“从\(X\)到它的第\(i\)个坐标集\(X_i\)的\DefineConcept{投射}”.
\end{definition}
